% !TEX root = ../main.tex

\section{Related Work}


% \subsection{From Similarity Metrics to Multi-dimensional Evaluation}

% "System outputs that are more similar to the references are deemed better." \citep{gehrmann2023repairing}

% Early methods for automated text evaluation centered on lexical overlap, most prominently the BiLingual Evaluation Understudy (BLEU) \citep{papineni_bleu_2001} and Recall-Oriented Understudy for Gisting Evaluation (ROUGE) \citep{lin_rouge_2004}. BLEU was developed for machine translation and evaluates candidate translations by measuring modified (n)-gram precision against human references, while ROUGE was developed for summarization and measures how much reference content is recovered by a candidate summary. Both metrics therefore operationalize quality through surface-level similarity to reference texts.

% However, the efficacy of metrics based on lexical overlap is limited. Surface-level textual features do not capture rich semantic information inherent to language, limiting their validity as measures of text quality \citep{reiter2018structured}. This motivated the development of embedding-based evaluation methods such as BERTScore \citep{zhang_bertscore_2020}, BLEURT \citep{sellam_bleurt_2020}, and MoverScore \citep{zhao_moverscore_2019}, which use contextual embeddings to estimate semantic similarity between candidate and reference texts.

% One major limitation of similarity-based metrics is that they may not capture all relevant information needed to assess the evaluation target. As \citet{gehrmannRepairingCrackedFoundation2022} note, quality differences between advanced NLG models are not reflected in surface-level textual features. This led to the development of more fine-grained evaluation methods which decompose the evaluation target into multiple criteria. In a seminal paper, \citet{fabbri_summeval_2021} introduce an elaborate benchmark for summarization models, where they decompose "summary quality" into the subdimensions \textit{coherence}, \textit{consistency}, \textit{fluency} and \textit{relevance}. They re-evaluate a wide array of summarization models on these dimensions by using expert annotators, exposing specific failure modes: pre-trained language models correlated highly with human ratings on consistency and fluency, while extractive models struggled with coherence and relevance, highlighting their known shortcomings. 

\subsection{From Similarity Metrics to Multi-dimensional Evaluation}

Early methods for automated text evaluation centered on lexical overlap, most prominently the BiLingual Evaluation Understudy (BLEU) \citep{papineni_bleu_2001} and Recall-Oriented Understudy for Gisting Evaluation (ROUGE) \citep{lin_rouge_2004}. Both metrics operationalize quality through surface-level similarity to reference texts: BLEU measures modified $n$-gram precision against human translations, while ROUGE measures how much reference content is recovered by a candidate summary.

However, the efficacy of lexical-overlap metrics is limited. Surface-level textual features do not capture the rich semantic information inherent to language, limiting their validity as measures of text quality \citep{reiter2018structured}. This motivated embedding-based evaluation methods such as BERTScore \citep{zhang_bertscore_2020}, BLEURT \citep{sellam_bleurt_2020}, and MoverScore \citep{zhao_moverscore_2019}, which use contextual embeddings to estimate semantic similarity between candidate and reference texts.

Yet, similarity-based metrics remain limited because holistic similarity scores may not capture all relevant information needed to assess the evaluation target. As \citet{gehrmannRepairingCrackedFoundation2022} note, quality differences between advanced NLG models are not fully reflected in surface-level textual features. This motivated more fine-grained evaluation methods that decompose the evaluation target into multiple criteria. A central example is SummEval, where \citet{fabbri_summeval_2021} decompose "summary quality" into \textit{coherence}, \textit{consistency}, \textit{fluency}, and \textit{relevance}, allowing them to identify distinct model failure modes.

\subsection{Automated Multi-dimensional Evaluators}

The discussion sparked by \citet{fabbri_summeval_2021} led to the development of automated multi-dimensional text evaluators. Early studies develop dimension-specific evaluators, such as for \textit{factual consistency} in summarization \citep{kryscinskiEvaluatingFactualConsistency2020, wangAskingAnsweringQuestions2020}, or \textit{coherence} in dialogue generation \citep{dziriEvaluatingCoherenceDialogue2019, yeQuantifiableDialogueCoherence2021}. Later studies worked towards unified measurement, such as \citet{zhongUnifiedMultiDimensionalEvaluator2022a}, who propose a method to quantify any arbitrary dimension using transformer models. 

Although their method correlates highly with expert annotations, it requires a costly pre-training stage and reference documents. The introduction of LLMs, known for their general problem-solving capabilities, offered a possible solution to these two problems. Early research into using LLMs as reference-free NLG evaluators \citep{fuGPTScoreEvaluateYou2023, wangChatGPTGoodNLG2023} exhibited lower correlations with expert judgments compared to neural evaluators such as \citet{zhongUnifiedMultiDimensionalEvaluator2022a}. However, with the release of GPT-4, \citet{liu_g-eval_2023} outperformed both reference-based and reference-free baselines by a substantial amount. \citet{zheng_judging_2023} introduce novel benchmarks to analyse NLG model performance on open-ended NLG tasks and show that LLM evaluations obtain an agreement with human preferences comparable to human-to-human agreement, solidifying the use of LLMs as general-purpose evaluators. 

\subsection{Textual Quantification in the Social Sciences}

The potential of LLMs for quantifying textual attributes is not limited to benchmarking NLG models. Given the function of language as a means for expression, linguistic analysis offers a window into human behavior. Text data analysis is used in psychology to study human behavioral traits, such as personality \citep{boydLanguagebasedPersonalityNew2017}, psychological well-being \citep{guntukuDetectingDepressionMental2017}, and emotional states \citep{kahnMeasuringEmotionalExpression2007}; in economics, to study stock market dynamics using text from news and social media \citep{manelaNewsImpliedVolatility2017}; and in political science, to study partisanship using congressional speeches \citep{gentzkowMeasuringGroupDifferences2019}. Given the abundance of text data and the costliness of manual annotation, NLP methods are therefore a natural choice to facilitate the measurement of complex constructs from text. 

The social sciences faced similar limitations as the NLG community. Dictionary methods, which consider basic features such as word count, are cheap to implement, but often too superficial to answer complex research questions. Representation-based methods, such as embedding models, are expensive to deploy, and most suitable for answering narrow research questions. LLMs again offered the potential for a one-size-fits-all solution: \citet{rathjeGPTEffectiveTool2024} show that GPT models can accurately measure numerous psychological constructs from texts in multiple different languages, and \citet{lemensUncoveringSemanticsConcepts2023} use GPT-4 to quantify semantical properties of texts, both studies achieving state-of-the-art performance. More generally, \citet{yangLLMMeasureGeneratingValid2024} propose a method to measure arbitrary concepts from text data leveraging the LLMs internal concept representation. 

\subsection{Conditions for Meaningful Measurement}

The current LLM-based methods for measuring constructs from text in the social sciences show many similarities to the evaluation of the capabilities of NLG models in ML: Both fields define measurement targets which they operationalize through LLM scoring. However, as a consequence of their measurement assumptions, they differ in the interpretations that the measurements permit.

The considerations in \citet{fabbri_summeval_2021} are mostly pragmatic. They observed that a single holistic similarity score did not sufficiently cover the quality differences between summarization systems, so they proposed to measure subdimensions that expose these differences. For evaluating models on these specific criteria, this is sufficient. However, their procedure falls short of permitting general claims about summary quality and summarization capabilities. 

As \citet{salaudeenMeasurementMeaningValidityCentered2025} argue, this inferential leap requires justifying the relationship between the measurements and the object of the claim; in this case, the relationship between criterion scores and the construct of "summarization ability". This practice is standard in the social sciences: \citet{boydLanguagebasedPersonalityNew2017, guntukuDetectingDepressionMental2017, kahnMeasuringEmotionalExpression2007, manelaNewsImpliedVolatility2017, gentzkowMeasuringGroupDifferences2019, rathjeGPTEffectiveTool2024, lemensUncoveringSemanticsConcepts2023, yangLLMMeasureGeneratingValid2024} all operate within a framework where the target constructs are theoretically defined, and where measurements are validated against established measures or theoretically relevant external criteria. Because of this, their measurements are not just criterion-level scores, but construct measurements whose interpretation is anchored in the broader theoretical frameworks of their respective fields.

\subsection{Meaningful Multi-Dimensional Measurement}\label{sec:litrev_meaning}

\paragraph{Criticisms of Multi-Dimensional LLM Evaluators}

To achieve the level of meaningful measurement seen in the social sciences, a growing number of studies are recognizing the need for validity-guided workflows for LLM evaluation based on measurement theory and psychometrics \citep{xiaoEvaluatingEvaluationMetrics2023, linPromptsConstructsDualValidity2025, binetteImprovingValidityPractical2024, beanMeasuringWhatMatters2025, salaudeenMeasurementMeaningValidityCentered2025, wangEvaluatingGeneralPurposeAI, kardanovaNovelPsychometricsBasedApproach2024, federiakinImprovingLLMLeaderboards2025, PDFModernLLM2026, feuerWhenJudgmentBecomes2025, alaaMedicalLargeLanguage2025}. Yet, advanced benchmarks do not immediately solve the problem of automated evaluation, particularly for open-ended language generation tasks, for which the set of possible correct answers is large and often unknown \citep{freitagExpertsErrorsContext2021}.

% Meanwhile, a growing number of studies have recognized the need for validity-guided workflows based on the principles of measurement theory and psychometrics for designing benchmarks for general-purpose AI models such as LLMs \citep{xiaoEvaluatingEvaluationMetrics2023, linPromptsConstructsDualValidity2025, binetteImprovingValidityPractical2024, beanMeasuringWhatMatters2025, salaudeenMeasurementMeaningValidityCentered2025, wangEvaluatingGeneralPurposeAI, kardanovaNovelPsychometricsBasedApproach2024, federiakinImprovingLLMLeaderboards2025, PDFModernLLM2026, feuerWhenJudgmentBecomes2025, alaaMedicalLargeLanguage2025}. Yet, advanced benchmarks do not immediately solve the problem of automated evaluation, particularly for open-ended language generation tasks, for which the set of possible correct answers is large and often unknown \citep{freitagExpertsErrorsContext2021}. 

As it stands, the focus in automated NLG evaluation still predominantly revolves around theory-agnostic multi-dimensional measurement protocols, where criteria are measured in isolation and the relationship between the measurement and the evaluation target is underspecified. As \citet{arias2025statistical} point out, this holds for 88\% of the papers published in NLP evaluation as of 2024, while only 7.6\% employ true multicriteria methods, even though linguistic quality is inherently multi-dimensional. This distinction marks the central problem for meaningful multi-dimensional measurement: decomposition is necessary, but not sufficient; additionally, the criteria must be related back to the target construct.

\paragraph{Relating Criteria to Constructs}

Importantly, this relationship defines the level at which the claims are made: at the system level, or at the artefact level. Benchmarks concern system-level claims, as they attempt to measure some latent model ability. Psychometric methods lend themselves well to this, as many have been developed for this specific purpose. For example, in the same way that latent writing ability can be inferred from grading an essay on multiple criteria, LLMs can be used as scorers to achieve the same goal \citep{wang2025evaluating}. However, these methods are not designed for artefact-level claims. Scores that are valid for comparing systems do not automatically support diagnostic claims about individual artefacts. For example, when a model obtains a high score on writing ability, it does not automatically imply that every single text it produces is of high quality, nor does it highlight what needs to be improved about the text. For diagnostic purposes, therefore, different methods are required.

The field has proposed several approaches, with varying measurement assumptions. For automated evaluation of a dialogue generation task, \citet{hashemi_llm-rubric_2024} propose a method to model rater idiosyncrasies using multi-dimensional LLM scores. They collect data from 24 humans who annotate an average of 31 dialogues on their overall quality. By using criterion-level LLM scores and fitting them to the holistic human ratings using a calibration network, they can accurately model individual rater behavior, enabling them to model for example expert annotators. However, they explicitly take the stance that they model subjective rating behavior, not some ground truth. \citet{kolaSAJASimpleApproach2026} take a similarly pragmatic approach, where they reuse the same six evaluation criteria for evaluation across multiple distinct evaluation tasks, achieving state-of-the-art alignment with human ratings. They argue that the six-dimensional LLM ratings contain all required information, and a lightweight calibration head is sufficient to extract it. 

\paragraph{Meaningful Artefact-Level Diagnostics}

\citet{hashemi_llm-rubric_2024} and \citet{kolaSAJASimpleApproach2026} are examples of the frontier of multi-dimensional evaluation for diagnostic purposes, and they confirm that decomposition has potential for improving LLM evaluators. However, the pragmatism of their methods limits their utility. First of all, they both use calibration heads as their recomposition function. Although this might improve alignment with human ratings, it lacks explainability. Because of this, although the accuracy of their single predicted holistic quality score is high, it is difficult to determine which dimensions contributed to this score, and how. Moreover, their approaches explicitly model human idiosyncrasies. This is not necessarily wrong; human annotations are often considered the gold standard in alignment research. However, it is also widely acknowledged that human labels are inherently noisy \citep{zhaoChallengesTrustworthyHuman}, with recent evidence suggesting that LLMs can sometimes outperform expert coders on difficult annotation tasks \citep{tornbergLargeLanguageModels2025}. Therefore, methods are needed that distinguish alignment with human ratings from meaningful measurement of the target construct, treating human judgments as one source of validity evidence rather than as the object of measurement itself.

\citet{bang_transforming_2025} take an important step in this direction with UniScore, which derives transparent criterion weights from their discriminative power against an external signal. However, while this makes recomposition interpretable, it leaves the measurement model implicit. Because the criteria remain user-defined, the framework can show how strongly each criterion relates to an external signal, but not whether the selected criteria adequately define the construct that the final score is interpreted as measuring. This is a central validity risk for artefact-level diagnostics: when an overall score is constructed from multiple criteria, the choice of criteria is itself part of the construct specification.

The present thesis addresses this problem by treating artefact-level LLM evaluation as a construct-specification problem rather than as just an aggregation problem. It examines whether theoretically grounded criteria, measured separately through LLM prompts and recomposed into an overall score, provide stronger validity evidence and more interpretable artefact-level measurements than holistic or weakly specified alternatives. To formalize this approach, artefact-level scoring is specified as a formative measurement model, drawing on established model specifications from psychometrics.
